We evaluate \condAgenticCogWriter\ around three questions: (1) does process-level organization improve long-form document quality relative to alternative writing systems, (2) which proposed components contribute to any advantage, and (3) how does the writing loop behave in practice? Seven systems were specified before scoring, and one additional execution variant was added after the main results as an exploratory diagnostic. Every comparison uses the same task inputs and information policy.

\subsection{Systems Compared}
\label{sec:baselines}

\paragraph{Alternative writing systems.}
We compare \condAgenticCogWriter\ with three ways of organizing long-form generation: (1) \condSinglePass\ produces the complete document in one generation call; (2) \condStaged\ executes \texttt{Pre-Write}$\rightarrow$\texttt{Write}$\rightarrow$\texttt{Re-Write}, following the common separation of planning, drafting, and revision~\citep{yang2022re3,wan2025cognitive}; and (3) \condTaskPlanning\ follows a WriteHERE-style adaptive task graph~\citep{xiong2025beyond}, recursively decomposing the assignment and selecting among dependency-satisfied task nodes. The alternatives therefore advance a whole document, a fixed writing stage, or an adaptive task node, whereas \condAgenticCogWriter\ advances a writing process.

\paragraph{Ablations of \condAgenticCogWriter.}
Two pre-specified ablations change one control mechanism while retaining the rest of the writing loop. \condNoGoals\ removes the explicit goal representation but keeps adaptive process selection and the three subagents. \condFixedOrder\ keeps the shared draft, goals, process history, and subagents but replaces the \texttt{Monitor}'s adaptive decision with a repeated \texttt{Planning}$\rightarrow$\texttt{Translating}$\rightarrow$\texttt{Reviewing} cycle. The comparison tests whether explicit goals or adaptive ordering explain the full system's advantage.

\paragraph{Execution variants.}
Two additional systems probe how the selected process is executed. \condSingleWriter\ retains the goal representation and adaptive process loop but performs planning, drafting, and reviewing in one writer context without subagent delegation; it also limits each \texttt{Translating} step to one paragraph, so the condition is diagnostic rather than a clean ablation. \condSingleContext\ was added after the main results as a narrower exploratory check: it preserves the full process loop and shared state but executes each selected process in the \texttt{Monitor}'s context rather than a separate subagent context. Because its run also differs in gateway entry point and transport, we use it only to test whether a simple role-context explanation remains plausible.

\begin{table*}[t]
\centering
\small
\setlength{\tabcolsep}{4pt}
\begin{tabular}{lllll}
\toprule
\textbf{Condition}
& \textbf{Control unit}
& \textbf{Goal network}
& \textbf{Scheduling}
& \textbf{Execution} \\
\midrule
\multicolumn{5}{l}{\textit{Baselines}} \\
\midrule
\condSinglePass
& Document
& \no
& \no
& One context \\
\condStaged~\cite{yang2022re3,wan2025cognitive}
& Stage
& \no
& Fixed stages
& One context \\
\condTaskPlanning~\cite{xiong2025beyond}
& Task node
& \no
& Adaptive tasks
& Delegated tasks \\
\midrule
\multicolumn{5}{l}{\textit{Process-level conditions}} \\
\midrule
\condFull
& Writing process
& \yes
& Adaptive processes
& Delegated roles \\
\condNoGoals
& Writing process
& \no
& Adaptive processes
& Delegated roles \\
\condFixedOrder
& Writing process
& \yes
& Fixed cycle
& Delegated roles \\
\midrule
\multicolumn{5}{l}{\textit{Diagnostic ablation}} \\
\midrule
\condSingleWriter
& Paragraph
& \yes
& Writer-decided
& One context \\
\bottomrule
\end{tabular}
\caption{Experimental conditions. The primary distinction is the controller's unit of action. The process-level interventions remove the explicit goal network or adaptive process scheduling while keeping delegated role execution. \condSingleWriter\ removes delegation but also changes the control unit, so it does not isolate process separation.}
\label{tab:experimental-conditions}
\end{table*}


\subsection{Benchmarks and Prompts}
\label{sec:benchmarks}

We use three benchmarks to test whether process-level organization transfers across different forms of long-form writing. WritingBench~\citep{wu2026writingbench} covers broad real-world writing requests with task-specific evaluation criteria; HelloBench~\citep{que2024hellobench} includes long-text summarization, open-ended generation, and completion; and DoLoMiTes~\citep{malaviya2025dolomites} targets methodical expert writing with explicit objectives, procedures, inputs, and constraints. Together, the benchmarks cover general instruction following, sustained long-context generation, and structured expert-document production.

We focus evaluation on tasks most likely to expose differences in how systems manage long documents. Before any system output was scored, we removed prompts used while developing the protocol and deterministically prioritized prompts with the longest supplied contexts and, when the task states a target length, the longest requested outputs. WritingBench therefore combines long inputs with explicit long-output requirements, HelloBench emphasizes very long-input summarization and continuation, and DoLoMiTes emphasizes procedural constraints with more moderate prompt length. Appendix~\ref{app:prompts} gives the exact selection rule, sample sizes, and sensitivity analysis.

\subsection{Implementation}
\label{sec:implementation}

All systems run in the same Codex-style coding-agent harness. We use \modelGenBare\ as the generator and \modelJudgeBare\ as the primary evaluator; \modelCross\ is used only for the cross-family robustness check.\footnote{The generator is held fixed across systems so that the experiment tests writing organization rather than base-model capability. We use the higher-capability evaluator for document assessment and separately check selected comparisons with a different model family~\citep{openai2026luna,openai2026sol,openai2026gpt56}.} \condAgenticCogWriter\ is packaged as an agent skill: the top-level \texttt{Monitor} invokes \texttt{Planner}, \texttt{Translator}, and \texttt{Reviewer} through the harness's subagent mechanism. The comparison systems use the same harness and task-level information policy. Appendix~\ref{app:prompts} reproduces the full agent-skill, subagent, baseline, and evaluator prompts; Appendix~\ref{app:runtime} reports the runtime settings.

Each prompt--system pair is generated in three runs with the same generator, prompt set, and runtime configuration. Repeating generation lets us measure variation due to stochastic execution while keeping the task set and system definition fixed. Network access, web search, and external retrieval are disabled for every system, and the harness records completed outputs plus available process traces, token accounting, and timing.

\subsection{Evaluation}
\label{sec:evaluation}

\paragraph{Pointwise evaluation.}
Pointwise evaluation scores each completed document independently. We report benchmark-native criteria where available and a common five-dimension rubric across all three benchmarks; the benchmark-specific procedures and complete evaluator prompts are in Appendix~\ref{app:prompts}. The common rubric scores instruction fulfillment, organization and coherence, content adequacy and depth, style and audience fit, and factual or constraint fidelity from 1 to 5. Because the raw ratings concentrate near the upper end of the scale, each dimension is standardized within a benchmark across outputs from all seven pre-specified systems before the five standardized scores are averaged. The resulting composite is relative to the compared systems within that benchmark rather than an absolute quality scale.

\paragraph{Pairwise evaluation.}
Pairwise evaluation directly compares responses written for the same prompt. We ask the evaluator to compare two complete responses using the same quality dimensions and return a winner or tie. LLM judges can exhibit position bias~\citep{zheng2023judging,wang2024fair}, so every pair is evaluated in both presentation orders. We retain a winner only when both orders select the same semantic response; any disagreement, including winner versus tie, is recorded as a tie. Headline win rates condition on comparisons with an agreed winner.

The primary generator and evaluator belong to the same model family, which can create systematic preference leakage or self-preference~\citep{li2026preference,panickssery2024selfpreference}. We therefore repeat selected system comparisons with \modelCross\ as a descriptive cross-family robustness check. The two judges are reported separately rather than averaged.

\paragraph{Pre-specified statistical testing.}
We fixed the confirmatory comparisons and tests before scoring to avoid selecting favorable analyses after observing results. Within each benchmark, every pre-specified system contrast is tested separately in each generation run with an exact two-sided sign test over non-tied prompt-level wins and losses~\citep{dixonmood1946}. Because testing several related contrasts increases the chance of false positives, we apply Holm's step-down multiple-comparison correction~\citep{holm1979} across the pre-specified contrast-by-run tests within each benchmark. Generation runs remain separate in confirmatory inference because repeated generations of the same prompts are not independent prompt samples. Appendix~\ref{app:pairwise-detail} reports generation-run estimates, Wilson intervals~\citep{wilson1927}, and multiplicity-adjusted decisions.

The pre-specified tests are intentionally conservative about repeated generations. To estimate the overall size of the two ablation effects, we additionally report a post-hoc prompt-level summary: for each prompt, win/tie/loss outcomes are mapped to $1$, $1/2$, and $0$, averaged across generation runs, and then combined across benchmarks. We report a prompt-level bootstrap interval and sign test. Because this analysis was chosen after the main results were available, we treat it as exploratory rather than confirmatory.

\paragraph{Robustness and behavioral analysis.}
Automatic pairwise judging can be sensitive to output length and judge identity~\citep{zheng2023judging,dubois2024length,li2026preference,panickssery2024selfpreference}. We therefore test whether the pairwise direction persists within pre-specified output-length bands, under a tight compute match, and under the cross-family evaluator. Output length and realized compute are consequences of execution, so the first two analyses are descriptive sensitivity checks rather than causal adjustments.

We analyze observable traces of process transitions, goal creation and development, goal replacement, and subagent invocations to characterize how the writing loop is actually used. We also compare runs that complete one planning--drafting--reviewing cycle with runs that take another process path. The trace-defined split is observational rather than randomized, so it describes behavior without identifying a causal mechanism.
